% 3-3-1-benchmark.tex
%  Budget: 3pp.  The primary dataset -- every frozen number rests on it.
%  Source map: the open benchmark dataset, its DE series and the two-year
%   held-out test protocol -> lago_forecasting_2021 ;
%   negative prices / volatility as intrinsic market features -> jiang_review_2018 .
%  Numbers verified 2026-09-05 against data/raw/DE.csv (the pinned snapshot):
%   52416 rows, 0 NaN, 0 duplicate stamps; train 34944h / 297 neg (0.8499%),
%   test 17472h / 241 neg (1.3793%); overall 538 neg (1.0264%);
%   mean 34.6779 std 15.9400 min -221.99 max 210.00 median 33.45.

\subsection{مجموعه‌داده بنچمارک \lr{EPEX-DE}}
\label{sec:3-3-1}

\subsubsection*{ساختار و دامنه}

داده‌ی اصلی این پژوهش، سری ساعتیِ قیمت بازار روز-پیشِ آلمان در بورس \lr{EPEX} است،
که در قالب مجموعه‌داده‌ی باز لاگو و همکاران و همراه با پروتکل ارزیابیِ آن منتشر شده
است \cite{lago_forecasting_2021}. انتخاب این مجموعه به‌جای گردآوریِ مستقلِ داده، یک
تصمیم روش‌شناختی است: تنها در این حالت است که اعداد فصل چهارم را می‌توان مستقیماً و
بدون هیچ ابهامِ تعریفی در کنار اعداد منتشرشده‌ی مرجع نشاند.

مجموعه از ۹ ژانویه‌ی ۲۰۱۲ تا ۳۱ دسامبر ۲۰۱۷ را پوشش می‌دهد و شامل \lr{52{,}416} سطر ساعتی
است. هر سطر سه متغیر دارد:

\begin{itemize}
  \item \lr{price} \lr{—} قیمت تسویه‌ی ساعتیِ بازار روز-پیش، به یورو بر مگاوات‌ساعت؛
        متغیر هدف پژوهش.
  \item \lr{exog\_1} \lr{—} پیش‌بینی روزپیشِ بار مصرفی برای ناحیه‌ی کنترلیِ
        \lr{Amprion}.
  \item \lr{exog\_2} \lr{—} پیش‌بینی روزپیشِ تولید بادی و خورشیدی.
\end{itemize}

هر دو متغیر برون‌زا \emph{پیش‌بینی}اند \lr{—} نه مقدار محقق‌شده \lr{—} و پیش از حراج روز-پیش
منتشر می‌شوند. این نکته صرفاً توصیفی نیست؛ مبنای مجاز بودنِ استفاده از مقادیر
\emph{خودِ روز هدف} این دو متغیر در ماتریس ویژگی است، که \cref{sec:3-4} آن را به‌تفصیل
توجیه می‌کند.

نامِ ناحیه‌ایِ متغیر نخست را باید به‌خاطر سپرد. \lr{exog\_1} بارِ ناحیه‌ی
\lr{Amprion} است و نه بارِ ملیِ آلمان \lr{—} تمایزی که در دوره‌ی بنچمارک هیچ اثری
ندارد، چون آموزش و آزمون هر دو روی همان تعریف انجام می‌شوند، اما در آزمون
برون‌توزیعیِ \cref{sec:5-2} به هسته‌ی اصلیِ یافته بدل می‌شود.

\subsubsection*{تقسیم آموزش و آزمون}

تقسیم داده مطابق پروتکل مرجع انجام می‌شود: دو سال پایانی به‌طور کامل به‌عنوان دوره‌ی
آزمون کنار گذاشته می‌شود. دوره‌ی آزمون از ۴ ژانویه‌ی ۲۰۱۶ تا ۳۱ دسامبر ۲۰۱۷ ادامه
دارد و شامل ۷۲۸ مبدأ پیش‌بینی و \lr{17{,}472} ساعت است؛ دوره‌ی آموزش \lr{34{,}944} ساعت (\lr{1{,}456} روز)
پیش از آن را در بر می‌گیرد.

این کنارگذاری کامل است. هیچ مشاهده‌ای از دوره‌ی آزمون در برازش مدل‌ها، در تنظیم
ابرپارامترها، در برازش وزن‌های مدل ترکیبی و در تعیین آستانه‌ی رژیم نقشی ندارد. پنجره‌ی
تنظیم ابرپارامتر (۵ ژانویه‌ی ۲۰۱۵ تا ۳ ژانویه‌ی ۲۰۱۶) و پنجره‌ی برازش وزن‌ها (۱۲
ژانویه‌ی ۲۰۱۵ تا ۳ ژانویه‌ی ۲۰۱۶) هر دو اکیداً پیش از آغاز دوره‌ی آزمون پایان می‌یابند،
و این قید در کد با یک تابع تصریحی اعمال می‌شود، نه صرفاً رعایت می‌گردد (\cref{sec:3-5}).

\subsubsection*{آمار توصیفی قیمت}

میانگین قیمت در کل دوره \lr{34.68} یورو بر مگاوات‌ساعت با انحراف معیار \lr{15.94} است، و
میانه \lr{33.45} \lr{—} یعنی توزیع در ناحیه‌ی مرکزی تقریباً متقارن است. اما دنباله‌های آن
چنین نیستند: کمینه منفیِ \lr{221.99} و بیشینه \lr{210.00} یورو بر مگاوات‌ساعت است، چولگی
منفیِ \lr{0.52} و کشیدگی \lr{16.27}. کشیدگیِ بسیار بالا در کنار چولگیِ منفی، تصویر روشنی
می‌دهد: سری بیشتر وقت‌ها آرام است و گاه با جهش‌های شدید \lr{—} به‌ویژه رو به پایین و
تا ناحیه‌ی منفی \lr{—} از آن آرامش خارج می‌شود.

قیمت‌های منفی در این بازار پدیده‌ای واقعی‌اند و نه خطای داده: ۵۳۸ ساعت از کل دوره
(\lr{1.03} درصد) قیمت منفی دارند، که ۲۹۷ ساعت آن در دوره‌ی آموزش و ۲۴۱ ساعت در دوره‌ی آزمون
قرار می‌گیرد. نسبت آن‌ها در دوره‌ی آزمون بیشتر است (\lr{1.38} در برابر \lr{0.85} درصد)، یعنی
پدیده در طول دوره تشدید می‌شود. ادبیات نیز همین ویژگی‌ها \lr{—} نوسان شدید،
جهش‌های ناگهانی و ناهمگنی واریانس \lr{—} را از خصوصیات ذاتی سری قیمت برق برمی‌شمارد
\cite{jiang_review_2018}.

وجود قیمت منفی یک پیامد مستقیم بر انتخاب معیار دارد: هر معیاری که بر خودِ قیمت تقسیم
کند، در همسایگی صفر تعریف‌ناپذیر یا انفجاری می‌شود. \cref{sec:3-5} بر همین پایه \lr{MAPE}
را کنار می‌گذارد.

سطح قیمت در طول دوره ثابت نمی‌ماند. میانگین سالانه از \lr{42.89} در سال ۲۰۱۲ به \lr{37.78} در
۲۰۱۳، \lr{32.76} در ۲۰۱۴، \lr{31.63} در ۲۰۱۵ و \lr{28.98} در ۲۰۱۶ کاهش می‌یابد و آنگاه در ۲۰۱۷ به
\lr{34.20} بازمی‌گردد. این روند نزولی و بازگشت آن، همان چیزی است که فرض ایستایی در صورت
قوی خود را رد می‌کند (\cref{sec:3-2}) و واسنجی مجدد روزانه را ضروری می‌سازد.

\subsubsection*{کیفیت و بازتولیدپذیری}

مجموعه‌ی داده کامل است: هیچ مقدار گم‌شده‌ای در هیچ‌یک از سه ستون وجود ندارد و هیچ
برچسب زمانی تکراری نیست. اسکریپت وارسیِ مستقلی همین را راستی‌آزمایی می‌کند که
پیش از هر اجرای نتایج، ابعاد، بازه‌ی زمانی و پیوستگیِ نمایه‌ی ساعتی را می‌سنجد.

بازتولیدپذیری در این پژوهش با یک قاعده‌ی سخت اداره می‌شود: داده‌ی بیرونی یک‌بار
دریافت و در یک تصویرِ ثابت با اثر انگشت \lr{sha256} ثبت می‌شود، در مخزن پروژه
نگهداری می‌گردد، و از آن پس تنها از همان‌جا خوانده می‌شود. علت این سخت‌گیری آن است
که یک واسط برنامه‌نویسیِ زنده فردا ممکن است پاسخ دیگری بدهد و این کار را در سکوت
انجام دهد؛ آنگاه نتایجِ منتشرشده دیگر بازتولیدپذیر نیستند بی‌آنکه کسی متوجه شود.
دریافت مجدد این مجموعه در ۲۹ اوت ۲۰۲۶ بایت‌به‌بایت با تصویر ثبت‌شده یکسان بود.

افزون بر آن، کد مهندسی ویژگی از دسترسی به شبکه منع شده است و این منع با آزمون خودکار
اعمال می‌شود، نه با قرارداد. توضیح این سازوکار در \cref{sec:3-4} آمده است.
